% LaTeX Document for VESPER Interaction and Time Tracking System
% For inclusion in research paper

\documentclass[conference]{IEEEtran}
\usepackage{amsmath,amssymb,amsfonts}
\usepackage{algorithmic}
\usepackage{graphicx}
\usepackage{textcomp}
\usepackage{xcolor}
\usepackage{booktabs}
\usepackage{multirow}
\usepackage{listings}

% For code listings
\lstset{
    basicstyle=\ttfamily\footnotesize,
    breaklines=true,
    columns=fullflexible,
    keepspaces=true,
}

\begin{document}

% ============================================================================
% SECTION: INTERACTION AND TIME TRACKING SYSTEM
% ============================================================================

\section{VESPER Interaction and Time Tracking System}
\label{sec:interaction_system}

To enhance the realism and evaluation capabilities of the VESPER smart home simulation, we implemented a comprehensive interaction and time tracking system that integrates CASAS-compatible item sensors, virtual smart home devices, and temporal acceleration mechanisms. This system enables the generation of standardized activity datasets comparable to real-world smart home deployments while maintaining temporal efficiency in simulation.

\subsection{System Architecture}

The interaction system comprises three integrated subsystems, as illustrated in Figure~\ref{fig:interaction_architecture}:

\begin{enumerate}
    \item \textbf{Item Sensor Manager}: Implements CASAS-compatible binary sensors for object interaction tracking
    \item \textbf{Virtual Device Manager}: Simulates IoT device control and automation
    \item \textbf{Virtual Time Manager}: Enables temporal acceleration while maintaining timestamp accuracy
\end{enumerate}

% Architecture could be illustrated with a figure
% \begin{figure}[t]
% \centering
% \includegraphics[width=\columnwidth]{figures/interaction_architecture.pdf}
% \caption{VESPER Interaction System Architecture}
% \label{fig:interaction_architecture}
% \end{figure}

\subsection{Item Sensor Implementation}

Following the CASAS smart home dataset format~\cite{cook2012casas}, we implemented 19 binary item sensors distributed across five room types (Table~\ref{tab:item_sensors}). Each sensor tracks object interactions with ON/OFF states, timestamps, and interaction durations.

\begin{table}[t]
\centering
\caption{CASAS-Compatible Item Sensor Configuration}
\label{tab:item_sensors}
\begin{tabular}{@{}llll@{}}
\toprule
\textbf{Sensor ID} & \textbf{Object} & \textbf{Room} & \textbf{Category} \\
\midrule
\multicolumn{4}{l}{\textit{Kitchen (7 sensors)}} \\
I001 & Sink & Kitchen & Appliance \\
I002 & Stove & Kitchen & Appliance \\
I003 & Refrigerator & Kitchen & Appliance \\
I004 & Microwave & Kitchen & Appliance \\
I005 & Coffee Maker & Kitchen & Appliance \\
I006 & Kettle & Kitchen & Item \\
I007 & Dishes & Kitchen & Item \\
\midrule
\multicolumn{4}{l}{\textit{Dining Room (2 sensors)}} \\
I008 & Phone & Dining Room & Item \\
I009 & Dining Table & Dining Room & Furniture \\
\midrule
\multicolumn{4}{l}{\textit{Bathroom (4 sensors)}} \\
I010 & Bathroom Sink & Bathroom & Appliance \\
I011 & Shower & Bathroom & Appliance \\
I012 & Toilet & Bathroom & Appliance \\
I013 & Medicine Cabinet & Bathroom & Item \\
\midrule
\multicolumn{4}{l}{\textit{Bedroom (3 sensors)}} \\
I014 & Bed & Bedroom & Furniture \\
I015 & Closet & Bedroom & Furniture \\
I016 & Lamp & Bedroom & Item \\
\midrule
\multicolumn{4}{l}{\textit{Living Room (3 sensors)}} \\
I017 & Television & Living Room & Appliance \\
I018 & Couch & Living Room & Furniture \\
I019 & Book & Living Room & Item \\
\bottomrule
\end{tabular}
\end{table}

\subsubsection{Proximity-Based Interaction Detection}

Item sensors activate when the virtual agent approaches within a configurable proximity threshold (1.0--1.5 meters, depending on object type). The system supports two interaction modes:

\begin{itemize}
    \item \textbf{Automatic}: Sensors activate/deactivate based on proximity alone (e.g., furniture, frequently-used items)
    \item \textbf{Manual}: Sensors require explicit VLM decision-making for activation (e.g., appliances, task-specific objects)
\end{itemize}

The interaction detection algorithm is formalized as:

\begin{equation}
    \text{Interact}(o, a) = \begin{cases}
        \text{True}, & \text{if } \|p_a - p_o\| \leq d_o \\
        \text{False}, & \text{otherwise}
    \end{cases}
\end{equation}

where $p_a$ is the agent position, $p_o$ is the object position, and $d_o$ is the object-specific interaction distance threshold.

\subsubsection{CASAS Data Format Export}

The system exports interaction logs in CASAS-compatible format:

\begin{lstlisting}
2024-10-14 12:00:00.123 I008 Phone ON
2024-10-14 12:05:15.456 I008 Phone OFF
2024-10-14 12:10:30.789 I001 Sink ON
2024-10-14 12:12:45.012 I001 Sink OFF
\end{lstlisting}

Each entry contains: timestamp (millisecond precision), sensor ID, sensor name, and state change. This format enables direct comparison with real CASAS smart home datasets and facilitates activity recognition model evaluation.

\subsection{Virtual Device Control}

To simulate smart home automation, we implemented 11 virtual IoT devices across five rooms (Table~\ref{tab:virtual_devices}). Devices automatically activate based on task context and room location.

\begin{table}[t]
\centering
\caption{Virtual Smart Home Device Configuration}
\label{tab:virtual_devices}
\begin{tabular}{@{}llll@{}}
\toprule
\textbf{Device ID} & \textbf{Name} & \textbf{Type} & \textbf{Room} \\
\midrule
D001 & Kitchen Light & Light & Kitchen \\
D002 & Kitchen Stove & Appliance & Kitchen \\
D003 & Kitchen Fridge & Appliance & Kitchen \\
D004 & Kitchen Microwave & Appliance & Kitchen \\
D005 & Living Light & Light & Living Room \\
D006 & Living TV & Appliance & Living Room \\
D007 & Living Lamp & Light & Living Room \\
D008 & Bedroom Light & Light & Bedroom \\
D009 & Bedroom Lamp & Light & Bedroom \\
D010 & Bathroom Light & Light & Bathroom \\
D011 & Dining Light & Light & Dining Room \\
\bottomrule
\end{tabular}
\end{table}

\subsubsection{Task-Based Automation}

Devices activate automatically based on task semantics and location. For example:

\begin{itemize}
    \item \textbf{``Cook oatmeal''}: Kitchen Light ON, Stove ON
    \item \textbf{``Make phone call''}: Dining Light ON
    \item \textbf{``Wash hands''}: Kitchen Light ON (if kitchen sink) or Bathroom Light ON
    \item \textbf{``Watch TV''}: Living Light ON, Living TV ON
\end{itemize}

The system maintains device state transitions, usage statistics, and exports SmartThings-compatible JSON logs for integration with IoT platforms.

\subsection{Virtual Time Management}

To enable efficient simulation of long-duration activities (e.g., sleeping 8 hours, cooking 45 minutes), we implemented a temporal acceleration system that compresses virtual time while maintaining accurate timestamps for sensor events.

\subsubsection{Time Acceleration Algorithm}

Given a task with virtual duration $t_v$ and maximum allowed real duration $t_r^{\max}$, the system calculates the required time scale:

\begin{equation}
    s = \frac{t_v}{t_r^{\max}}
\end{equation}

where $s$ is the acceleration factor. For example, to simulate 8 hours of sleep in 5 seconds:

\begin{equation}
    s = \frac{28800 \text{ sec}}{5 \text{ sec}} = 5760\times
\end{equation}

The virtual timestamp at real time $t_r$ is computed as:

\begin{equation}
    t_v(t_r) = t_0 + s \cdot (t_r - t_{r0})
\end{equation}

where $t_0$ is the virtual start time and $t_{r0}$ is the real start time.

\subsubsection{Task-Specific Time Profiles}

We defined 24 task-specific time profiles based on typical ADL durations (Table~\ref{tab:time_profiles}). These profiles enable realistic temporal modeling while maintaining simulation efficiency.

\begin{table}[t]
\centering
\caption{Virtual Time Profiles for Common ADL Tasks}
\label{tab:time_profiles}
\begin{tabular}{@{}lrrr@{}}
\toprule
\textbf{Task Category} & \textbf{Virtual} & \textbf{Real} & \textbf{Scale} \\
& \textbf{Duration} & \textbf{Duration} & \textbf{Factor} \\
\midrule
\multicolumn{4}{l}{\textit{Sleep \& Rest}} \\
Sleep (8 hrs) & 28800s & 5s & 5760$\times$ \\
Nap (30 min) & 1800s & 3s & 600$\times$ \\
Rest (10 min) & 600s & 3s & 200$\times$ \\
\midrule
\multicolumn{4}{l}{\textit{Personal Care}} \\
Shower (10 min) & 600s & 4s & 150$\times$ \\
Wash hands (1 min) & 60s & 2s & 30$\times$ \\
Brush teeth (2 min) & 120s & 2s & 60$\times$ \\
\midrule
\multicolumn{4}{l}{\textit{Cooking}} \\
Cook simple (15 min) & 900s & 4s & 225$\times$ \\
Cook complex (45 min) & 2700s & 6s & 450$\times$ \\
Cook oatmeal (10 min) & 600s & 3s & 200$\times$ \\
Microwave (3 min) & 180s & 2s & 90$\times$ \\
\midrule
\multicolumn{4}{l}{\textit{Eating}} \\
Eat meal (20 min) & 1200s & 3s & 400$\times$ \\
Snack (5 min) & 300s & 2s & 150$\times$ \\
\midrule
\multicolumn{4}{l}{\textit{Communication}} \\
Phone call (5 min) & 300s & 3s & 100$\times$ \\
Long call (30 min) & 1800s & 5s & 360$\times$ \\
\midrule
\multicolumn{4}{l}{\textit{Entertainment}} \\
Watch TV (1 hr) & 3600s & 8s & 450$\times$ \\
Read (30 min) & 1800s & 4s & 450$\times$ \\
\midrule
\multicolumn{4}{l}{\textit{Cleaning}} \\
Clean dishes (5 min) & 300s & 2s & 150$\times$ \\
Clean room (15 min) & 900s & 4s & 225$\times$ \\
\bottomrule
\end{tabular}
\end{table}

\subsection{System Integration}

The interaction system integrates with the existing VLM-based navigation framework through six key integration points:

\begin{enumerate}
    \item \textbf{System Initialization}: Initialize item sensors, virtual devices, and time manager at BGE startup
    \item \textbf{Task Start}: Activate relevant sensors and devices when a new task begins
    \item \textbf{State Update}: Check proximity and trigger interactions during each navigation step
    \item \textbf{Task Completion (Success)}: Finalize interactions, log durations, and deactivate devices
    \item \textbf{Task Completion (Failure)}: Clean up partial interactions for failed tasks
    \item \textbf{Data Export}: Export all sensor logs, device states, and time records to CASAS-compatible formats
\end{enumerate}

\subsection{Data Export Formats}

The system generates four standardized output files:

\subsubsection{Item Sensor Log (CASAS Format)}
Text file with millisecond-precision timestamps:
\begin{lstlisting}[basicstyle=\ttfamily\scriptsize]
item_sensor_log_20241014_120000.txt
2024-10-14 12:00:00.123 I008 Phone ON
2024-10-14 12:05:15.456 I008 Phone OFF
\end{lstlisting}

\subsubsection{Detailed Interaction Log (JSON)}
Structured JSON with complete interaction metadata:
\begin{lstlisting}[language=json,basicstyle=\ttfamily\scriptsize]
{
  "session_start": "2024-10-14T12:00:00",
  "interactions": [
    {
      "sensor_id": "I008",
      "sensor_name": "Phone",
      "room": "dining_room",
      "start_time": "2024-10-14T12:00:00.123",
      "end_time": "2024-10-14T12:05:15.456",
      "duration_seconds": 315.333,
      "task": "Make a phone call",
      "interaction_type": "automatic"
    }
  ]
}
\end{lstlisting}

\subsubsection{Device Control Log (SmartThings Format)}
JSON log of all device state transitions:
\begin{lstlisting}[language=json,basicstyle=\ttfamily\scriptsize]
{
  "device_events": [
    {
      "device_id": "D011",
      "device_name": "Dining_Light",
      "action": "ON",
      "timestamp": "2024-10-14T12:00:00",
      "trigger": "task_start:Make a phone call"
    }
  ],
  "device_usage": {
    "D011": {"on_time": 315, "activations": 1}
  }
}
\end{lstlisting}

\subsubsection{Virtual Time Log}
Records time acceleration events and temporal mapping:
\begin{lstlisting}[language=json,basicstyle=\ttfamily\scriptsize]
{
  "time_acceleration_events": [
    {
      "task": "Make a phone call",
      "time_scale": 100.0,
      "real_duration": 3.0,
      "virtual_duration": 300.0,
      "start_real": "2024-10-14T12:00:00",
      "start_virtual": "2024-10-14T12:00:00",
      "end_virtual": "2024-10-14T12:05:00"
    }
  ]
}
\end{lstlisting}

\subsection{Evaluation Benefits}

This comprehensive interaction and time tracking system provides several advantages for ADL recognition research:

\begin{enumerate}
    \item \textbf{CASAS Compatibility}: Direct comparison with real smart home datasets
    \item \textbf{Temporal Realism}: Accurate representation of activity durations
    \item \textbf{Simulation Efficiency}: Complex multi-hour scenarios execute in seconds
    \item \textbf{Ground Truth Data}: Precise interaction timestamps and durations for evaluation
    \item \textbf{IoT Integration}: SmartThings-compatible device logs for automation research
    \item \textbf{Scalability}: Configurable sensors and devices for diverse scenarios
\end{enumerate}

The system enables rapid generation of large-scale, temporally-accurate ADL datasets that would require weeks or months to collect in real smart home deployments, while maintaining the data quality and format compatibility necessary for rigorous activity recognition evaluation.

% ============================================================================
% EXPERIMENTAL VALIDATION
% ============================================================================

\subsection{System Validation}

We validated the interaction system through a series of CASAS-aligned ADL tasks. Table~\ref{tab:validation_results} presents sample execution metrics for five common activities.

\begin{table}[t]
\centering
\caption{Interaction System Validation Results}
\label{tab:validation_results}
\begin{tabular}{@{}lrrrr@{}}
\toprule
\textbf{Task} & \textbf{Virtual} & \textbf{Real} & \textbf{Sensors} & \textbf{Devices} \\
& \textbf{Time} & \textbf{Time} & \textbf{Triggered} & \textbf{Used} \\
\midrule
Make phone call & 5m & 3s & 1 & 1 \\
Wash hands & 1m & 2s & 1 & 1 \\
Cook oatmeal & 10m & 3s & 3 & 2 \\
Eat meal & 20m & 3s & 1 & 1 \\
Clean dishes & 5m & 2s & 2 & 1 \\
\midrule
\textbf{Total} & 41m & 13s & 8 & 6 \\
\bottomrule
\end{tabular}
\end{table}

The results demonstrate that the system successfully compresses 41 minutes of virtual activity into 13 seconds of real execution time (189$\times$ average acceleration) while maintaining accurate sensor triggering and device control. All generated logs conform to CASAS format specifications and can be directly processed by existing activity recognition pipelines.

% ============================================================================
% REFERENCES (if needed in standalone section)
% ============================================================================

% \begin{thebibliography}{1}
% \bibitem{cook2012casas}
% D. J. Cook, A. S. Crandall, B. L. Thomas, and N. C. Krishnan, 
% ``CASAS: A smart home in a box,'' 
% \textit{IEEE Computer}, vol. 45, no. 7, pp. 62--69, 2012.
% \end{thebibliography}

\end{document}
